\documentclass[11pt]{article}
\usepackage{fontspec}
\usepackage[french]{babel}
\usepackage{amsmath}
\usepackage{amssymb}
\usepackage{amsthm}
\usepackage{mathrsfs}
\usepackage{tikz-cd}
\usepackage{enumitem}

\newcommand{\printedpage}[1]{\marginpar{\footnotesize p.~#1}}
\newcommand{\msp}[1]{\marginpar{\footnotesize [#1]}}
\newcommand{\hra}{\lhook\joinrel\longrightarrow}
\newcommand{\Cat}{\mathrm{Cat}}

\begin{document}
\setcounter{footnote}{76}

\printedpage{73}
% head: A. Grothendieck : Dérivateurs, Ch. XII, Édition des Universités de Montpellier II et Paris VII
% foot: Version 7 décembre 2002

\msp{page 105}

Je prendrai encore $(X,U,\{t\})$ tel que $t$ soit objet final de $X$, mais si on veut $q\notin W$, on ne peut supposer $t'$ élément final de $X'$. Au lieu de cela, je vais supposer

\begin{enumerate}[label=\alph*)]
\item $U$, $U'$ (et aussi $X$, ayant objet final $t$) $W$-asphériques, donc $i\in W$ (et même $i$ $W$-asphérique), $p\in W$, mais
\item $X'$ non $W$-asphérique, donc $i',q\notin W$.
\end{enumerate}

Pour ceci, notons que
\[ U'/t'\simeq U'/p_!(e_U)\overset{\text{déf}}{=}\underline{\pi}_0(U/U'). \]
Le corollaire de la proposition 7 (cf. p. 85), non encore démontré, nous dit que $i'\in W$ (i.e. ici, $X'$ $W$-asphérique) si et seulement si $i'$ est $W$-asphérique, i.e. $U'/t'$ est $W$-asphérique. Nous allons l'admettre ici. Tout revient donc à trouver une flèche
\[ p:U\longrightarrow U' \]
dans Cat avec les propriétés suivantes.

\begin{enumerate}[label=\alph*)]
\item $U$ et $U'$ sont $W$-asphériques,
\item $\underline{\pi}_0(U/U')$ n'est pas $W$-asphérique.
\end{enumerate}

Nous allons reprendre l'exemple précédent
\[ f:X\longrightarrow Y \]
(p. 101-103), en prenant
\[ U=X,\qquad \underline{\pi}_0(U/U')=Y, \]

\msp{page 106}

donc tout revient, à partir de l'exemple donné $X\xrightarrow{f}Y$, à construire une flèche
\[ \pi:Y\longrightarrow U' \tag*{$(*)$} \]
dans Cat qui satisfasse les conditions

\begin{enumerate}[label=\alph*)]
\item $\pi$ est un isomorphisme local,
\item $U'$ est $W$-asphérique.
\end{enumerate}

On prendra alors
\[ p=\pi f. \]
Mais par la façon dont $X$ a été construit à partir d'un $Y$ 1-connexe et non $\infty$-connexe, par ailleurs quelconque, on est ramené finalement à trouver $(*)$ satisfaisant a), b), et de plus

\begin{enumerate}[label=\alph*), start=3]
\item \printedpage{74}
% head: A. Grothendieck : Dérivateurs, Ch. XII, Édition des Universités de Montpellier II et Paris VII
% foot: Version 7 décembre 2002
$Y$ est 1-connexe, mais non $\infty$-connexe.
\end{enumerate}

(NB On suppose maintenant $W=W_\infty$.) Il revient au même de trouver

% printed with a vertical rule at the left of this block
\begin{quote}
$U'$ objet \textit{asphérique} de Cat, $F\in\operatorname{Ob}U'^\wedge$ tels que $U'/F$ soit 1-connexe et non $\infty$-connexe.
\end{quote}

On songe par exemple à prendre simplement une sous-catégorie ouverte $U'_0$ de $U'$, qui soit 1-connexe et non $\infty$-connexe. Partons d'une telle catégorie $U'_0$ - par exemple une réalisation homotopique dans Cat, de la sphère $S^2$. On l'immerge comme ouvert dans un $U'$ asphérique, en lui ajoutant un objet final. On a gagné.

\msp{page 107}

Faisons un énoncé en forme.

\textbf{Proposition 11.} \textit{On peut trouver une situation} $(*)$ (p. 124 \texttt{[plutôt 104]}) \textit{avec}

\begin{enumerate}[label=\textit{\alph*)}]
\item \textit{$t$ objet final de $X$.}
\item \textit{$U$, $U'$ (et $X$) $W_\infty$-asphériques, donc $i,p\in W_\infty$ (et même $i$ $W_\infty$-asphérique).}
\item \textit{$X'$ non $W_\infty$-asphérique, donc $i',q\notin W_\infty$.}
\end{enumerate}

Il suffit pour cela de partir d'un objet 1-connexe $Y$ de Cat, de lui adjoindre un objet final $u'_0$, ce qui donne $U'$, de prendre $U=\underline{\mathrm{Ch}}(Y;y_0,-)$, $p:U\longrightarrow U'$ le composé $U\xrightarrow{\text{but}}Y\underset{\text{can.}}{\overset{\text{incl.}}{\hra}}U'$, et $X$ déduit de $U$ en lui ajoutant un objet final $t$. On a alors $U'/t'\simeq Y$, d'où
\[ X'\ W_\infty\text{-asphérique (i.e. }i'\in W_\infty)\Longleftrightarrow Y\ W_\infty\text{-asphérique.} \]
On a le contre-exemple cherché en prenant $Y$ (1-connexe mais) non $\infty$-connexe.

\subsection*{2.4\quad Application aux immersions ouvertes dwyériennes.}

Donnons quand même l'enoncé \textit{positif} auquel nous en avions. % [sic] l'enoncé -> l'énoncé, missing accent

\textbf{Théorème 3.} \textit{Soient $(X,U,T)$ un triple $\in\mathcal{T}$ ($U$ ouvert de $X$ dans} Cat\textit{, $T=X\backslash U$), et pour toute flèche $p:U\longrightarrow U'$ flèche de} Cat\textit{, le carré cocartésien dans} Cat

\[
\begin{tikzcd}
U \arrow[r, hook, "i"] \arrow[d, "p"'] & X \arrow[d, "q"'] & T \arrow[l, dotted, hook'] \arrow[d, dotted, "\text{iso}"] \\
U' \arrow[r, hook, "i'"] & X' & T'. \arrow[l, dotted, hook']
\end{tikzcd}
\]

\printedpage{75}
% head: A. Grothendieck : Dérivateurs, Ch. XII, Édition des Universités de Montpellier II et Paris VII
% foot: Version 7 décembre 2002

\msp{page 108}

\textit{Supposons que pour tout $t\in\dot{T}=\dot{U}$, le préfaisceau d'ensembles}
\[ i^*(t)=(u\longmapsto\operatorname{Hom}(u,t))\in U^\wedge \]
\textit{soit totalement $W$-asphérique (déf. 1, p. 88)} (\footnote{\textbf{NB} On dit sous ces conditions que $i_0:U\longrightarrow\bar{U}$ est \textit{strictement $W$-asphérique} (cf. déf. 2 page suivante).}) \textit{- p. ex. que la catégorie $U/t$ soit totalement $W$-asphérique (prop. 8, page 93). (NB le cas où les $U/t$ ont un objet final, cf. prop. 9 a), i.e. celui où l'inclusion $U\longrightarrow\bar{U}$ admet un adjoint à droite, est celui envisagé par Dwyer, celui des immersions ouvertes de Dwyer} (\footnote{Je ne connais pas d'autre cas où th. 3 s'applique - faute de connaître de catégorie totalement $W$-asphérique qui n'ait un objet final, ni d'objet $F$ d'une catégorie $U^\wedge$ qui soit totalement $W$-asphérique sans être représentable, i.e. $U/F$ a un objet final. À présent, je sais que les catégories \textit{filtrantes} sont totalement $W$-asphériques si $W$ satisfait l'axiome des limites inductives.})\textit{.) Alors pour toute $p:U\longrightarrow U'$, l'immersion ouverte}
\[ i'_0:U'\longrightarrow\overline{U'} \]
\textit{est $W$-asphérique (et même, tout comme $i_0$, strictement $W$-asphérique), et a fortiori elle est} \texttt{[dans]} \textit{$W$. Par suite (prop. 6} \texttt{[p. 81]}\textit{) on a ceci.}

\[
\begin{aligned}
&\text{\textcircled{a}}\qquad && i\in W && \Longleftrightarrow\quad && i'\in W\quad && \textit{moyennant } \mathrm{W}(7\text{ bis}).\\
&\text{\textcircled{b}}\qquad && p\in W && \Longleftrightarrow\quad && q\in W\quad && \textit{moyennant } \mathrm{W}(7).\ (\footnotemark)
\end{aligned}
\]
\footnotetext{Démonstration page 112.}

\textbf{Corollaire 1.} \textit{Supposons} W(7)\textit{. Si $i:U\hra X$ est une immersion ouverte telle que $i_0:U\hra\bar{U}$ soit totalement} \texttt{[il suffit strictement]} \textit{$W$-asphérique, alors $i:U\longrightarrow X$ est une $W$-cofibration (`ultraforte'), par quoi j'entends ceci : pour toute flèche $U'\xrightarrow{f}U''$ dans $U\backslash$}Cat\textit{,}

\msp{page 109}

\textit{telle que $f\in W$, la flèche $g:X'\longrightarrow X''$ qui s'en déduit par cochangement de base à partir de l'objet $X$ de $U\backslash$}Cat\textit{, est aussi dans $W$.}

\textsc{Démonstration du corollaire} : on applique le théorème 3 à l'immersion ouverte $U'\overset{i'}{\hra}X'$ déduite de $U\overset{i}{\hra}X$ par le cochangement de base $U\xrightarrow{f}U'$, en notant de plus que l'hypothèse du théorème 3 sur $i$ est préservée en passant de $i$ à $i'$.

Il faut quand même introduire une terminologie adéquate :

\textbf{Définition 2 :}

\begin{enumerate}[label=\alph*)]
\item Une flèche $f:X\longrightarrow Y$ est dite \textit{strictement $W$-asphérique} si pour toute $y\in\operatorname{Ob}Y$, $f^*(y)=(x\longmapsto\operatorname{Hom}_Y(f(x),y))$ est un objet totalement $W$-asphérique de $X^\wedge$ (cf. déf. 1, page 88). Elle est dite \textit{totalement $W$-asphérique} si pour tout $y\in\operatorname{Ob}Y$, $X/y$ ($\simeq X/f^*(y)$) est un objet totalement $W$-asphérique de Cat.

\item \printedpage{76}
% head: A. Grothendieck : Dérivateurs, Ch. XII, Édition des Universités de Montpellier II et Paris VII
% foot: Version 7 décembre 2002
Une immersion ouverte $i:U\longrightarrow X$ est dite \textit{$W$-dwyérienne} (resp. \textit{totalement} \texttt{[$W$-]} \textit{dwyérienne}) si l'immersion induite $i_0:U\longrightarrow\bar{U}$ est strictement (resp. totalement) $W$-asphérique. Une immersion fermée $j:T\longrightarrow X$ est dite \textit{$W$-dwyérienne} (resp. \textit{totalement $W$-dwyérienne}) si et seulement si $j^\circ:T^\circ\longrightarrow X^\circ$ est une immersion ouverte $W$-dwyérienne (resp totalement $W$-dwyérienne).

\item Une flèche $i:U\longrightarrow X$ de Cat (pas nécessairement une immersion)

\msp{page 110}

est appelée une \textit{$W$-cofibration,} si pour toute flèche $f:U'\longrightarrow U''$ dans $U\backslash$Cat, la flèche $g:X'\longrightarrow X''$ qui s'en déduit par le cochangement de base par $X$ (objet de $U\backslash$Cat) est dans $W$ quand $f$ l'est.
\end{enumerate}

\textbf{NB} Cette notion de $W$-cofibration est stable par composition et par cochangement de base, et les isomorphismes de Cat sont des $W$-cofibrations. Si $f$ est une $W$-cofibration, $f^\circ$ l'est et inversement.

(Le corollaire s'énonce en disant qu'une immersion ouverte ou fermée qui est $W$-dwyérienne est une $W$-cofibration.)

\begin{enumerate}[label=\alph*), start=4]
\item Une flèche $i:U\longrightarrow X$ dans Cat est appelée une \textit{$W$-cofibration triviale}, ou une \textit{$W$-équivalence couniverselle,} si pour toute $U'$ dans $U\backslash$Cat, i.e. toute flèche $p:U\longrightarrow U'$ dans Cat, la flèche déduite de $i$,
\[ i':U'\longrightarrow X'=U'\amalg_U X'\quad\texttt{[plutôt }U'\amalg_U X\texttt{]}\ \ , \]
est dans $W$. (NB Il est immédiat que cela implique bien que $i$ est une $W$-cofibration.)
\end{enumerate}

(Cette notion de $W$-cofibration triviale est elle aussi stable par composition et par cochangement de base, et les isomorphismes dans Cat en sont.)

\msp{page 111}

\textbf{Corollaire 2.} (Moyennant W(7).) \textit{Soit $i:U\longrightarrow X$ une immersion ouverte $W$-dwyérienne. Pour que $i$ soit une $W$-cofibration triviale, il faut et il suffit que $i\in W$, ou ce qui revient au même} (cf. page 82), \textit{que $\dot{T}\hra T$ soit dans $W$.}

C'est évidemment nécessaire, et la suffisance résulte du théorème 3 a).

\textbf{Questions ouvertes.}

\begin{enumerate}[label=\alph*)]
\item J'ignore si une immersion ouverte $i$ qui est une $W$-cofibration, voire une $W$-cofibration triviale, est forcément $W$-dwyérienne (\footnote{et peut-être même une immersion de Dwyer?}). Mais par le corollaire de la proposition 7 (p. 85), je sais que c'est le cas si on suppose $T$ isomorphe à la catégorie finale $e$. Il s'ensuit toute au moins ceci : Pour que non seulement $i:U\longrightarrow X$, mais aussi toutes les immersions ouvertes $i_1:U\longrightarrow X_1$ qui s'en déduisent par le procédé de prop. 7 b) au moyen de flèches quelconques $T_1\longrightarrow T$ de but $T$ dans Cat, soient des $W$-cofibrations, il faut et il suffit que $i$ soit dwyérienne. % [sic] toute au moins -> tout au moins, agreement error

\item \printedpage{77}
% head: A. Grothendieck : Dérivateurs, Ch. XII, Édition des Universités de Montpellier II et Paris VII
% foot: Version 7 décembre 2002
Y-a-t-il une bonne notion d' immersion dwyérienne, applicable à toutes les immersions

\msp{page 112}

et incluant le cas des immersions ouvertes et des immersions fermées?
\end{enumerate}

\textsc{Démonstration} du théorème 3 (p. 107). Il suffit de prouver que si $i$ est $W$-dwyérienne (i.e. $i_0$ strictement $W$-asphérique). alors $i'$ l'est également - car il en résultera que $i'$ est $W$-asphérique et a fortiori $\in W$. Il faut donc vérifier que pour tout $t'\in\dot{U}'=\dot{T}$, l'élément $i^*(t')=(u'\longmapsto\operatorname{Hom}_{X'}(u',t'))$ de $U'^\wedge$ est totalement $W$-asphérique. Or $t'$ provient d'un unique élément $t$ de $\dot{T}$, et on sait qu'on a % [sic] asphérique). alors -> asphérique), alors, period before lower-case word

\[ i'^*(t)=p_!(i^*(t))\quad\texttt{[plutôt }i'^*(t')=p_!(i^*(t))\texttt{]},\qquad\text{où }\ p_!:U^\wedge\longrightarrow U'^\wedge. \]

Par hypothèse, $F=i^*(t)$ est un élément totalement $W$-asphérique de $U^\wedge$, ce qui signifie que pour toute flèche $p:U\longrightarrow U'$ dans Cat de souce $U$, $p_!(F)$ est un élément $W$-asphérique de $U'$. En fait, on a le lemme immédiat (par transitivité $(p'p)_!\simeq p'_!p_!$). % [sic] souce -> source, missing letter

\textbf{Lemme.} \textit{Soit $U$ dans} Cat\textit{, et $F$ dans $U^\wedge$. Si $F$ est totalement $W$-asphérique (resp. si $U/F$ est totalement $W$-asphérique), alors pour toute flèche $p:U\longrightarrow U'$ dans} Cat \textit{de source $U$, $F'=p_!(F)$ est un objet totalement $W$-asphérique de $U'^\wedge$ (resp.}

\msp{page 113}

\textit{$U'/F'$ est un objet totalement $W$-asphérique de} Cat\textit{.)}

\textbf{Corollaire 3} (du théorème 3). \textit{Si $i:U\longrightarrow X$ est une immersion ouverte ou fermée $W$-dwyérienne (resp. totalement $W$-dwyérienne), alors pour toute flèche $p:U\longrightarrow U'$ dans} Cat\textit{, $i':U'\longrightarrow X'=U'\amalg_U X'$} \texttt{[plutôt $U'\amalg_U X$]} \textit{l'est aussi.}

(NB Le cas non respé est essentiellement le théorème 3. Il ne dépend d'aucun axiome supplémentaire sur $W$, mais n'a d'intérêt qu'en vue des équivalences \textcircled{a} et \textcircled{b} (p. 108) qui en résultent, et qui sont valables moyennant W(7 bis) resp. W(7) seulement.)

\textbf{Commentaire de motivation :} les immersions ouvertes ou fermées $W$-dwyériennes, ou totalement $W$-dwyériennes, sont stables par cochangement de base (d'après le corollaire précédent), mais n'ont pas du tout l'air de l'être par composition. À cet égard, les `immersions de Dwyer' (où on suppose l'existence de l'adjoint à droite de $i_0:U\longrightarrow\bar{U}$, dans le cas d'une immersion ouverte) se comportent mieux, étant stables par composition. (\textbf{NB} Il ne semble pas

\msp{page 114}

non plus clair que les flèches strictement $W$-asphériques, ou les totalement $W$-asphériques, soient stables par composition - sorry!) Ma motivation pour m'échiner sur la notion $W$-dwyérienne est double.

\begin{enumerate}[label=\alph*)]
\item J'aimerais avoir des conditions à la fois compréhensibles, et aussi générales que possible, sur une immersion (ouverte p. ex), pour que ce soit une $W$-cofibration, propriété visiblement importante (toute comme la notion duale de $W$-fibration). % [sic] toute comme -> tout comme, agreement error
\end{enumerate}

\end{document}
